\chapter{Cenários e critérios de aceitação}
\label{cap:casos}

\section{Cenários}
\label{sec:cenarios}

Os cenários simulados, no presente trabalho, representam configurações do
ambiente a simular que, depois de iniciadas, são observadas pelos sensores.
Apenas a informação efetivamente detetada é posteriormente transformada em
eventos e factos utilizados pelo motor de inferência. Neste sentido, existe
uma diferença entre aquilo que existe no ambiente e aquilo que o sistema
conhece num determinado instante, ou seja, o carregamento de um cenário não
significa que o motor de inferência tenha imediatamente acesso a todos os
elementos existentes no mesmo. Por exemplo, um sinal ou um peão pode estar
presente no percurso, mas apenas passa a influenciar a tomada de decisão
quando é detetado pelos sensores e a respetiva informação chega ao motor de
inferência.

Na versão atual do projeto encontram-se definidos cerca de catorze cenários,
de diferente natureza. Apresentam-se na \cref{tab:cenarios} alguns exemplos.

\begin{table}[H]
  \centering
  \caption{Exemplos de cenários e respetivo comportamento}
  \label{tab:cenarios}
  \small
  \begin{tabularx}{\textwidth}{@{}L{5.2cm}VX@{}}
    \toprule
    \textbf{Cenário} & \textbf{Comportamento} \\
    \midrule
    Estrada sem situações de perigo & Continuação da circulação na ausência de
      condições que exijam intervenção \\ \linhasep
    Alteração do limite de velocidade & Deteção do novo limite e adaptação da
      velocidade \\ \linhasep
    Semáforo vermelho & Reconhecimento da obrigação de parar \\ \linhasep
    Sinal STOP & Aproximação, paragem e continuação após o cumprimento da
      condição \\ \linhasep
    Peão ou objeto na via & Avaliação do risco e escolha da ação em função da
      situação \\ \linhasep
    TTC crítico & Reação perante uma situação de elevado risco de colisão \\ \linhasep
    Mudança de via impedida & Avaliação da ocupação da via adjacente \\ \linhasep
    Chuva & Adaptação do comportamento às condições meteorológicas \\ \linhasep
    Cenário sequencial & Resposta do sistema a diferentes situações
      introduzidas ao longo do mesmo percurso \\
    \bottomrule
  \end{tabularx}
\end{table}

Para além de definir os elementos inicialmente presentes no ambiente, os
cenários com ações temporizadas permitem introduzir alterações durante o
decorrer da simulação, como modificar as condições meteorológicas ou
adicionar elementos ao percurso. Desta forma, o ambiente não permanece
necessariamente igual desde o início até ao fim do cenário, o que permite
observar a reação do sistema à evolução das condições. A utilização destes
cenários permite avaliar não só se uma determinada regra é ativada, mas
também se o sistema acompanha corretamente a evolução da situação. Num
cenário com um peão, por exemplo, interessa acompanhar o processo desde o
momento em que é detetado pelos sensores até à decisão. Num cenário de TTC
crítico, interessa também perceber de que forma as diferentes leituras
recebidas ao longo do tempo influenciam a resposta do sistema.

\section{Critérios de aceitação}
\label{sec:criterios}

Os cenários que incluem várias situações ao longo do mesmo percurso permitem
perceber a interação entre os diferentes tipos de conhecimento e verificar se
o sistema consegue adaptar sucessivamente as suas decisões à medida que
surgem novos elementos no ambiente. No entanto, para além dos cenários, é
também importante ter em conta os critérios de aceitação do comportamento do
sistema, uma vez que não é suficiente verificar apenas se uma determinada
regra foi selecionada. É necessário confirmar se a informação utilizada pelo
motor de inferência corresponde ao estado observado, se a decisão sugerida é
coerente com as regras definidas e se o resultado pode ser explicado a partir
do conhecimento representado. Neste sentido, foram considerados os seguintes
critérios de aceitação:
\begin{itemize}
  \item a informação detetada pelos sensores deve dar origem a eventos e
    factos coerentes com a situação simulada;
  \item as regras a aplicar devem corresponder às condições presentes no
    cenário;
  \item a decisão final deve respeitar o nível de risco e as prioridades
    definidas para o sistema;
  \item quando existirem várias propostas de decisão, o mecanismo de
    resolução de conflitos deve selecionar a alternativa esperada;
  \item o comportamento do sistema deve acompanhar a evolução do cenário,
    nomeadamente as alterações que ocorram ao longo do tempo;
  \item após as decisões propostas deve ser permitido identificar, sempre que
    possível, a regra e a evidência que lhes deram origem;
  \item situações equivalentes devem dar origem a comportamentos consistentes
    em execuções repetidas;
  \item os cenários devem permitir identificar regras incompletas, conflitos
    inesperados ou condições ainda não representadas;
  \item à medida que a implementação em Prolog for concluída, deverá ser
    verificado se, perante factos conceptualmente equivalentes, as conclusões
    obtidas são coerentes com as obtidas no Drools.
\end{itemize}

\section{Testes realizados}
\label{sec:testes}

Para além dos cenários utilizados durante a simulação, foram realizados
testes para validar e verificar o comportamento das regras e a integração
entre os componentes. Atualmente, o projeto conta com cerca de 41 testes no
\emph{backend}, nos quais se incluem testes específicos às regras do Drools,
aos conflitos entre propostas de decisão, à descrição dos factos e testes
\emph{end-to-end} do motor de simulação e das mudanças de via. No
\emph{frontend} encontram-se ainda definidos 16 testes unitários, com foco nas
componentes e serviços da interface, bem como 7 testes \emph{end-to-end} com
Playwright, para testar situações como a introdução dos limites de velocidade
e o desvio de obstáculos. Os testes utilizados, até ao momento, permitem
verificar não só as regras de forma isolada, mas também a alteração do
ambiente, a inferência, a decisão e a respetiva representação na aplicação.
